\documentclass{amsart}
% For dealing with references we use the comment environment
\usepackage{verbatim}
\newenvironment{reference}{\comment}{\endcomment}
%\newenvironment{reference}{}{}
\newenvironment{slogan}{\comment}{\endcomment}
\newenvironment{history}{\comment}{\endcomment}

% For commutative diagrams we use Xy-pic
\usepackage[all]{xy}

% We use 2cell for 2-commutative diagrams.
\xyoption{2cell}
\UseAllTwocells

% We use multicol for the list of chapters between chapters
\usepackage{multicol}

% This is generally recommended for better output
\usepackage{lmodern}
\usepackage[T1]{fontenc}

% For cross-file-references

% Package for hypertext links:
\usepackage{hyperref}

% For any local file, say "hello.tex" you want to link to please

% Theorem environments.
%
\theoremstyle{plain}
\newtheorem{theorem}[subsection]{Theorem}
\newtheorem{proposition}[subsection]{Proposition}
\newtheorem{lemma}[subsection]{Lemma}

\theoremstyle{definition}
\newtheorem{definition}[subsection]{Definition}
\newtheorem{example}[subsection]{Example}
\newtheorem{exercise}[subsection]{Exercise}
\newtheorem{situation}[subsection]{Situation}

\theoremstyle{remark}
\newtheorem{remark}[subsection]{Remark}
\newtheorem{remarks}[subsection]{Remarks}

\numberwithin{equation}{subsection}

% Macros
%
\def\lim{\mathop{\mathrm{lim}}\nolimits}
\def\colim{\mathop{\mathrm{colim}}\nolimits}
\def\Spec{\mathop{\mathrm{Spec}}}
\def\Hom{\mathop{\mathrm{Hom}}\nolimits}
\def\Ext{\mathop{\mathrm{Ext}}\nolimits}
\def\SheafHom{\mathop{\mathcal{H}\!\mathit{om}}\nolimits}
\def\SheafExt{\mathop{\mathcal{E}\!\mathit{xt}}\nolimits}
\def\Sch{\mathit{Sch}}
\def\Mor{\mathop{\mathrm{Mor}}\nolimits}
\def\Ob{\mathop{\mathrm{Ob}}\nolimits}
\def\Sh{\mathop{\mathit{Sh}}\nolimits}
\def\NL{\mathop{N\!L}\nolimits}
\def\CH{\mathop{\mathrm{CH}}\nolimits}
\def\proetale{{pro\text{-}\acute{e}tale}}
\def\etale{{\acute{e}tale}}
\def\QCoh{\mathit{QCoh}}
\def\Ker{\mathop{\mathrm{Ker}}}
\def\Im{\mathop{\mathrm{Im}}}
\def\Coker{\mathop{\mathrm{Coker}}}
\def\Coim{\mathop{\mathrm{Coim}}}

% Boxtimes
%
\DeclareMathSymbol{\boxtimes}{\mathbin}{AMSa}{"02}

%
% Macros for moduli stacks/spaces
%
\def\QCohstack{\mathcal{QC}\!\mathit{oh}}
\def\Cohstack{\mathcal{C}\!\mathit{oh}}
\def\Spacesstack{\mathcal{S}\!\mathit{paces}}
\def\Quotfunctor{\mathrm{Quot}}
\def\Hilbfunctor{\mathrm{Hilb}}
\def\Curvesstack{\mathcal{C}\!\mathit{urves}}
\def\Polarizedstack{\mathcal{P}\!\mathit{olarized}}
\def\Complexesstack{\mathcal{C}\!\mathit{omplexes}}
% \Pic is the operator that assigns to X its picard group, usage \Pic(X)
% \Picardstack_{X/B} denotes the Picard stack of X over B
% \Picardfunctor_{X/B} denotes the Picard functor of X over B
\def\Pic{\mathop{\mathrm{Pic}}\nolimits}
\def\Picardstack{\mathcal{P}\!\mathit{ic}}
\def\Picardfunctor{\mathrm{Pic}}
\def\Deformationcategory{\mathcal{D}\!\mathit{ef}}

\setlength{\emergencystretch}{3em}
\begin{document}
\title[Stacks Project, Tag 0E7J]{493 --- Proper one-dimensional morphisms with fibrewise vanishing first cohomology and surjective structure-sheaf pushforward retain these properties universally near the fibre}
\maketitle
\noindent Copyright (C) 2005--2025 Johan de Jong. Stacks Project authors. Source: \href{https://stacks.math.columbia.edu/tag/0E7J}{Tag 0E7J}.

\noindent Editorial difficulty note (not author text). Difficulty H3: The proof must establish finite generation of the pushforward over a possibly non-Noetherian and infinitely generated target before applying Nakayama. An ambient polynomial embedding, Tor-independent derived base change and Noetherian approximation supply pseudo-coherence and finiteness, allowing fibrewise surjectivity to spread with universal cohomology vanishing..

\noindent Editorial provenance note (not author text). History: excerpt from revision \texttt{a0444\allowbreak 6e57e\allowbreak c1fbc\allowbreak 252a8\allowbreak 71afc\allowbreak ec775\allowbreak 2fb28\allowbreak 07b14}, more-morphisms.tex, lines 21058--21190. Reference presentation was adapted; original-author context and core support blocks were retained or added. The mathematical wording is unchanged. Licensed under GNU FDL 1.2 or later; no invariant sections or cover texts. See ../licenses/COPYING. Context blocks reproduce author definitions and supporting results; they are not additional counted problems.

\begin{lemma}
\label{lemma-h1-fibre-zero-h0-kappa}
Consider a commutative diagram
$$
\xymatrix{
X \ar[rr]_f \ar[rd] & & Y \ar[ld] \\
& S
}
$$
of morphisms of schemes. Let $s \in S$ be a point. Assume
\begin{enumerate}
\item $X \to S$ is locally of finite presentation and flat at
points of $X_s$,
\item $f$ is proper,
\item the fibres of $f_s : X_s \to Y_s$ have dimension $\leq 1$
and $R^1f_{s, *}\mathcal{O}_{X_s} = 0$,
\item $\mathcal{O}_{Y_s} \to f_{s, *}\mathcal{O}_{X_s}$ is surjective.
\end{enumerate}
Then there is an open $Y_s \subset V \subset Y$ such that
(a) $f^{-1}(V)$ is flat over $S$,
(b) $\dim(X_y) \leq 1$ for $y \in V$,
(c) $R^1f_*\mathcal{O}_X|_V = 0$,
(d) $\mathcal{O}_V \to f_*\mathcal{O}_X|_V$
is surjective,
and (b), (c), and (d) remain true after base change by any $Y' \to V$.
\end{lemma}

\begin{proof}
Let $y \in Y$ be a point over $s$.
It suffices to find an open neighbourhood of $y$
with the desired properties. As a first step, we replace $Y$
by the open $V$ found in Lemma \href{https://stacks.math.columbia.edu/tag/0E7G}{lemma h1 fibre zero (Tag 0E7G)} so that
$R^1f_*\mathcal{O}_X$ is zero universally (the hypothesis of
the lemma holds by Lemma \href{https://stacks.math.columbia.edu/tag/0E7F}{lemma check h1 fibre zero (Tag 0E7F)}).
We also shrink $Y$ so that all fibres of $f$ have dimension $\leq 1$
(use Morphisms, Lemma
\href{https://stacks.math.columbia.edu/tag/02FZ}{lemma openness bounded dimension fibres (Tag 02FZ)}
and properness of $f$). Thus we may assume we have (b) and (c)
with $V = Y$ and after any base change $Y' \to Y$.
Thus by Lemma \href{https://stacks.math.columbia.edu/tag/0E7I}{lemma h1 fibre zero check h0 kappa (Tag 0E7I)}
it now suffices to show (d) over $Y$.
We may still shrink $Y$ further; for example, we may and do
assume $Y$ and $S$ are affine.

\medskip\noindent
By Theorem \href{https://stacks.math.columbia.edu/tag/0399}{theorem openness flatness (Tag 0399)}
there is an open subset $U \subset X$ where $X \to S$
is flat which contains $X_s$ by hypothesis.
Then $f(X \setminus U)$ is a closed subset
not containing $y$. Thus after shrinking $Y$
we may assume $X$ is flat over $S$.

\medskip\noindent
Say $S = \Spec(R)$. Choose a closed immersion $Y \to Y'$ where $Y'$
is the spectrum of a polynomial ring $R[x_e; e \in E]$ on a set $E$.
Denote $f' : X \to Y'$ the composition of $f$ with $Y \to Y'$.
Then the hypotheses (1) -- (4) as well as (b) and (c)
hold for $f'$ and $s$. If we we show $\mathcal{O}_{Y'} \to f'_*\mathcal{O}_X$
is surjective in an open neighbourhood of $y$, then the same is true for
$\mathcal{O}_Y \to f_*\mathcal{O}_X$. Thus we may assume $Y$ is
the spectrum of $R[x_e; e \in E]$.

\medskip\noindent
At this point $X$ and $Y$ are flat over $S$. Then $Y_s$ and $X$
are tor independent over $Y$. We urge the reader to find their own
proof, but it also follows from
Lemma \href{https://stacks.math.columbia.edu/tag/0CTA}{lemma case of tor independence (Tag 0CTA)} applied to the
square with corners $X, Y, S, S$ and its base change by $s \to S$.
Hence
$$
Rf_{s, *}\mathcal{O}_{X_s} = L(Y_s \to Y)^*Rf_*\mathcal{O}_X
$$
by Derived Categories of Schemes, Lemma
\href{https://stacks.math.columbia.edu/tag/08IB}{lemma compare base change (Tag 08IB)}.
Because of the vanishing already established this implies
$f_{s, *}\mathcal{O}_{X_s} = (Y_s \to Y)^*f_*\mathcal{O}_X$.
We conclude that $\mathcal{O}_Y \to f_*\mathcal{O}_X$ is a map of
quasi-coherent $\mathcal{O}_Y$-modules whose pullback
to $Y_s$ is surjective. We claim
$f_*\mathcal{O}_X$ is a finite type $\mathcal{O}_Y$-module.
If true, then the cokernel $\mathcal{F}$ of
$\mathcal{O}_Y \to f_*\mathcal{O}_X$
is a finite type quasi-coherent $\mathcal{O}_Y$-module
such that $\mathcal{F}_y \otimes \kappa(y) = 0$.
By Nakayama's lemma (Algebra, Lemma \href{https://stacks.math.columbia.edu/tag/00DV}{lemma NAK (Tag 00DV)})
we have $\mathcal{F}_y = 0$. Thus $\mathcal{F}$ is zero
in an open neighbourhood of $y$
(Modules, Lemma \href{https://stacks.math.columbia.edu/tag/01B9}{lemma finite type stalk zero (Tag 01B9)})
and the proof is complete.

\medskip\noindent
Proof of the claim.
For a finite subset $E' \subset E$ set $Y' = \Spec(R[x_e; e \in E'])$.
For large enough $E'$ the morphism $f' : X \to Y \to Y'$ is proper, see
Limits, Lemma \href{https://stacks.math.columbia.edu/tag/0EX1}{lemma finite type eventually proper (Tag 0EX1)}.
We fix $E'$ and $Y'$ in the following.
Write $R = \colim R_i$ as the colimit of its finite type
$\mathbf{Z}$-subalgebras. Set $S_i = \Spec(R_i)$
and $Y'_i = \Spec(R_i[x_e; e \in E'])$.
For $i$ large enough we can find a diagram
$$
\xymatrix{
X \ar[d] \ar[r]_{f'} & Y' \ar[d] \ar[r] & S \ar[d] \\
X_i \ar[r]^{f'_i} & Y'_i \ar[r] & S_i
}
$$
with cartesian squares such that $X_i$ is flat over $S_i$ and
$X_i \to Y'_i$ is proper. See
Limits, Lemmas \href{https://stacks.math.columbia.edu/tag/01ZM}{lemma descend finite presentation (Tag 01ZM)},
\href{https://stacks.math.columbia.edu/tag/04AI}{lemma descend flat finite presentation (Tag 04AI)}, and
\href{https://stacks.math.columbia.edu/tag/081F}{lemma eventually proper (Tag 081F)}.
The same argument as above shows $Y'$ and $X_i$
are tor independent over $Y'_i$ and hence
$$
R\Gamma(X, \mathcal{O}_X) =
R\Gamma(X_i, \mathcal{O}_{X_i})
\otimes^\mathbf{L}_{R_i[x_e; e \in E']}
R[x_e; e \in E']
$$
by the same reference as above. By Cohomology of Schemes, Lemma
\href{https://stacks.math.columbia.edu/tag/02O6}{lemma proper over affine cohomology finite (Tag 02O6)}
the complex $R\Gamma(X_i, \mathcal{O}_{X_i})$ is pseudo-coherent
in the derived category of the Noetherian ring $R_i[x_e; e \in E']$ (see
More on Algebra, Lemma \href{https://stacks.math.columbia.edu/tag/066E}{lemma Noetherian pseudo coherent (Tag 066E)}).
Hence $R\Gamma(X, \mathcal{O}_X)$ is pseudo-coherent in the derived
category of $R[x_e; e \in E']$, see
More on Algebra, Lemma \href{https://stacks.math.columbia.edu/tag/0650}{lemma pull pseudo coherent (Tag 0650)}.
Since the only nonvanishing cohomology module
is $H^0(X, \mathcal{O}_X)$ we conclude it is a finite
$R[x_e; e \in E']$-module, see
More on Algebra, Lemma \href{https://stacks.math.columbia.edu/tag/064T}{lemma n pseudo module (Tag 064T)}.
This concludes the proof.
\end{proof}
\end{document}
